\chapter{Desenvolupament del programa}

Un cop definida la metodologia, vaig passar a la part pràctica del treball: el desenvolupament del programa en Python. Aquesta secció descriu com està construït el codi, quines llibreries he fet servir i com es processen les dades fins a obtenir el resultat final.

\section{Estructura general}
El programa segueix una estructura senzilla però efectiva:
\begin{enumerate}
    \item Connectar-se a la font de dades oberta (API de CityBikes).
    \item Rebre les dades en format JSON.
    \item Recórrer totes les estacions i obtenir el nombre de bicicletes i d'espais lliures.
    \item Calcular el nombre de canvis ràpids possibles a cada estació.
    \item Mostrar els resultats de manera clara i entenedora per pantalla.
\end{enumerate}

\section{Llibreries utilitzades}
He utilitzat principalment la llibreria \texttt{requests}, que permet fer peticions HTTP a servidors web i obtenir-ne les respostes. També he fet servir la llibreria \texttt{json}, que ve inclosa amb Python i serveix per treballar amb aquest format de dades.

Aquestes llibreries són molt utilitzades en el món de la programació i m’han permès simplificar molt la feina.

\section{Connexió amb l’API}
Per obtenir les dades he fet una crida a la URL següent:

\begin{verbatim}
url = "https://api.citybik.es/v2/networks/bicing"
response = requests.get(url)
\end{verbatim}

Aquesta línia estableix la connexió i descarrega el contingut. Tot seguit, comprovo que la resposta sigui correcta i converteixo les dades en un objecte JSON:

\begin{verbatim}
if response.status_code != 200:
    print("Error en la connexió:", response.status_code)
    exit()

data = response.json()
\end{verbatim}

\section{Recorregut de les estacions}
Les estacions estan guardades dins de \texttt{data["network"]["stations"]}. Per cada estació puc consultar valors com el nom, les bicicletes lliures i els espais disponibles.

\begin{verbatim}
for station in data["network"]["stations"]:
    bicis = station["free_bikes"]
    espais = station["empty_slots"]
\end{verbatim}

\section{Càlcul del canvi ràpid}
El canvi ràpid s’ha definit com el mínim entre les bicicletes lliures i els espais disponibles. Això s’implementa amb la funció \texttt{min} de Python:

\begin{verbatim}
canvi_rapid = min(bicis, espais)
\end{verbatim}

D’aquesta manera, si una estació té 5 bicis i 7 espais, els canvis ràpids possibles són 5. Si té 0 bicis i 12 espais, el resultat és 0.

\section{Sortida de resultats}
Per fer més entenedora la informació, el programa mostra per pantalla les dades amb un format clar:

\begin{verbatim}
print("Estació:", station["name"])
print("   -> Bicis lliures:", bicis)
print("   -> Espais lliures:", espais)
print("   -> Canvis ràpids possibles:", canvi_rapid)
print("-" * 40)
\end{verbatim}

La sortida és fàcil de llegir i permet veure ràpidament l’estat de cada estació.

\section{Exemple de funcionament}
Quan s’executa el programa, es pot obtenir una sortida similar a la següent:

\begin{verbatim}
Estació: Gran Via - Marina
   -> Bicis lliures: 5
   -> Espais lliures: 7
   -> Canvis ràpids possibles: 5
----------------------------------------
Estació: Passeig de Gràcia - Casp
   -> Bicis lliures: 0
   -> Espais lliures: 12
   -> Canvis ràpids possibles: 0
----------------------------------------
\end{verbatim}

Aquest exemple mostra com el programa processa correctament la informació i calcula els canvis ràpids a partir de les dades reals proporcionades per l’API.
